\section{Changing the child exponent with the height}
\label{sec:variable}
\label{var:section}

We now change the child aggregation to an $\ell_p$ norm, choosing
$p=1+1/h$ on a tree of height $h$.  This choice keeps the powers
$p^h\le e$ that arise in iteration bounded independently of $h$.
The resulting outer Hilbert sum is reflexive and admits no equivalent
AUC norm. Its midpoint estimates will be uniform in the component
height. We first consider centers with finite total support and
perturbations that vanish on a finite subtree containing the center's support in
each nonzero component. A normalized boundary calculation bounds the
sixth power of the perturbation by the excess in the sum of the squared
endpoint norms. Accumulating the energy that has left these subtrees
gives cubic control of the radius of a ball containing both endpoints.
We then obtain a bound for arbitrary centers with weakly null
perturbations, and a positive averaged midpoint modulus.

\subsection{The norm and its completion}

For an integer $h\ge1$, let
\[
 T_h=\bigcup_{j=0}^{h}\mathbb N^j,\qquad p_h=1+\frac1h.
\]
The root $o$ is the empty sequence; the children of $s$ are the
sequences $si$, $i\in\mathbb N$, when $|s|<h$.  We write $t\preceq s$
when $t$ is a prefix of $s$, including $t=s$.  All vectors and arrays
in this section are real.  For $v\in\ell_{p_h}(T_h)$, define two
nonnegative arrays from the leaves upward.  At a node $s$, put
\[
 a_s=\biggl(\sum_i P_{si}(v)^{p_h}\biggr)^{1/p_h},\qquad
 b_s=\biggl(\sum_i Q_{si}(v)^{p_h}\biggr)^{1/p_h},
\]
where both values are zero at a leaf, and set
\begin{equation}\label{var:recursion}
 P_s(v)=\max\{a_s,v_s+b_s\},\qquad
 Q_s(v)=\max\{a_s-v_s,b_s\}.
\end{equation}
Thus $P_s(v)-Q_s(v)=v_s$.  The pair is the coordinatewise least
nonnegative pair with this difference and with each array dominating
the $\ell_{p_h}$ norm of its children.  Indeed, at a node whose child
values have been fixed, every feasible pair must have
$P_s\ge\max\{a_s,v_s+b_s\}$ and $Q_s=P_s-v_s$.
The formulas are increasing in both child bounds, so induction from
the leaves proves minimality throughout the tree.  At least one of
the two child inequalities is an equality at every node.

Define
\[
 \nu_h(v)=P_o(v)+Q_o(v),\qquad
 E_h=(\ell_{p_h}(T_h),\nu_h),\qquad
 X_{\mathrm v}=\Bigl(\bigoplus_{h=1}^{\infty}E_h\Bigr)_{\ell_2}.
\]
We write $\|x\|_{\mathrm v}^2=\sum_h\nu_h(x^{(h)})^2$ for its norm.

\begin{proposition}\label{var:space}
The recursion \eqref{var:recursion} defines a norm equivalent to the
usual $\ell_{p_h}(T_h)$ norm.  In particular, $E_h$ is complete and
reflexive, with separable dual.  The space $X_{\mathrm v}$ is complete
and reflexive, its dual is separable, and vectors with finite total
coordinate support are dense in it.  Every finite coordinate
projection is bounded and has finite rank.
\end{proposition}

\begin{proof}
Write $p=p_h$ and $M_s=\max\{P_s(v),Q_s(v)\}$.  The recursion gives
\[
 M_s\le |v_s|+\|(M_{si})_i\|_p.
\]
Taking the $\ell_p$ norm on a level and proceeding upward proves
that all the defining arrays have finite level norms and that
\[
 M_o\le\sum_{j=0}^{h}\|(v_s)_{|s|=j}\|_p.
\]
In particular, the recursion is defined for every $v\in\ell_p(T_h)$.
Conversely, iteration of the child inequalities gives, on every
level $j$,
\[
 \|(P_s(v))_{|s|=j}\|_p\le P_o(v),\qquad
 \|(Q_s(v))_{|s|=j}\|_p\le Q_o(v).
\]
Since $v=P(v)-Q(v)$, these estimates imply
\begin{equation}\label{var:equivalence}
 (h+1)^{-1/p}\|v\|_{\ell_p(T_h)}
 \le\nu_h(v)
 \le 2(h+1)^{1-1/p}\|v\|_{\ell_p(T_h)}.
\end{equation}
They also give $|v_s|\le\nu_h(v)$ for every coordinate $s$.

Negation exchanges the two arrays, and positive scaling scales them.
If $(P,Q)$ and $(P',Q')$ are feasible pairs for $v$ and $w$, then
$(P+P',Q+Q')$ is feasible for $v+w$, by the triangle inequality in
$\ell_p$.  Minimality therefore gives the triangle inequality for
$\nu_h$.  The lower bound in \eqref{var:equivalence} gives
definiteness.  Norm equivalence proves completeness of $E_h$ and
density of its finitely supported vectors.

On a countable set, the usual dual of $\ell_p$, $1<p<\infty$,
is $\ell_{p/(p-1)}$: H\"older gives the pairing in one direction,
and testing a functional on finitely supported signed powers of its
coordinate values gives the converse and the dual norm.  Applying
this twice proves reflexivity, and finite support approximation
proves dual separability.  Equivalent norms preserve the continuous
dual, bidual, and surjectivity of the canonical map into the bidual,
so these properties hold for $E_h$.

The outer-sum argument in Proposition~\ref{rec:completion}, applied
to these component norms, proves completeness and reflexivity of
$X_{\mathrm v}$ and the dual identification
\[
 X_{\mathrm v}^*=\Bigl(\bigoplus_{h=1}^{\infty}E_h^*\Bigr)_{\ell_2}
\]
under the componentwise pairing. Finite component approximation and
separability of each $E_h^*$ give separability of this dual.
Truncating the component sum and then approximating within the retained
components proves finite total support density in $X_{\mathrm v}$.
Coordinate evaluations are bounded by the norm; each finite coordinate
projection is consequently a finite sum of bounded rank-one operators.
\end{proof}

The spaces also retain uniformly bounded paths with weakly null
increments.  This supplies their obstruction to asymptotically
uniformly convex renormings before we turn to their midpoint gains.

\begin{proposition}\label{var:no-auc}
The space $X_{\mathrm v}$ admits no equivalent asymptotically uniformly
convex norm.  More precisely, if
$a\|\cdot\|_{\mathrm v}\le N\le b\|\cdot\|_{\mathrm v}$ with
$0<a\le b$, then $\overline\delta_N(a/b)=0$.
\end{proposition}

\begin{proof}
For $s\in T_h$, let $e_s$ be its coordinate vector and put
$b_s=\sum_{t\preceq s}e_t$.  The path indicator itself is a feasible
$P$ array for $b_s$, with $Q=0$, and its root coordinate supplies the
reverse inequality. Hence $\nu_h(b_s)=1$. For a nonterminal $s$,
the vector $e_{si}$ has fields $P=1,Q=0$ at $si$ and $P=Q=1$ at
each proper ancestor. Thus $\nu_h(e_{si})=2$. The sequence of
distinct child vectors $(e_{si})_i$ is weakly null in $E_h$, by
\eqref{var:equivalence} and $p_h>1$, and hence in $X_{\mathrm v}$.
Since $b_{si}-b_s=e_{si}$ and the component heights are unbounded,
Lemma~\ref{rec:bounded-tree} applies with $m=M=1$, $c=2$,
$\alpha=a$, and $\beta=b$. It gives
$\overline\delta_N(a/b)=0$, excluding an AUC norm.
\end{proof}

\subsection{A boundary quantity and a sixth-power estimate}

Fix $h$ and abbreviate $p=p_h$.  A finite set $A\subset T_h$ is
\emph{ancestral} if it contains the root and every prefix of each of
its elements.  In the next two arguments, the center $x\in E_h$ is
supported on $A$, while $y\in E_h$ vanishes on $A$.  Write
\[
 B=\nu_h(x),\quad Y=\nu_h(y),\quad
 M_\sigma=\nu_h(x+\sigma y)\quad(\sigma\in\{+,-\}).
\]
The arrays for $x$ vanish off $A$.  On $A$, monotonicity of the
recursion shows that the arrays for $x+\sigma y$ have the form
\begin{equation}\label{var:common-increments}
 P_s(x+\sigma y)=P_s(x)+d_s^\sigma,\qquad
 Q_s(x+\sigma y)=Q_s(x)+d_s^\sigma,
 \qquad d_s^\sigma\ge0.
\end{equation}
The increments are equal because the coordinate difference is still
$x_s$.  In particular, $M_\sigma=B+2d_o^\sigma\ge B$.

Let $H$ consist of the nodes outside $A$ whose parents belong to $A$.
This is a possibly infinite antichain.  At each node $s\in A$, fix
one of the labels $P,Q$ whose old value equals its child norm.
For $i\in H$, let $v_i^\sigma$ be the value at $i$ of the endpoint
array with the label chosen by its parent.  The same label is used
for both signs.  On the subtree below $i$, the center is zero, so
negation interchanges the two arrays and
\begin{equation}\label{var:boundary-signs}
 v_i^++v_i^-=P_i(y)+Q_i(y).
\end{equation}
We shall propagate one of these two boundary vectors to the root.

\begin{lemma}\label{var:sixth-component}
Let $x,y\in E_h$, with $x$ supported on a finite ancestral set $A$
and $y$ zero on $A$.  Put $B=\nu_h(x)$, $Y=\nu_h(y)$, and
$M_\sigma=\nu_h(x+\sigma y)$.  With $\mathrm e=\exp(1)$ and
\[
 c_6=2\exp(-(\mathrm e-1))4^{-6},\qquad
 U^2=M_+^2+M_-^2,
\]
one has
\begin{equation}\label{var:sixth-component-eq}
 U^2-2B^2\ge c_6 U^{-4}Y^6.
\end{equation}
A quotient with numerator and denominator both zero is interpreted
as zero.  The same estimate holds when $x=0$ and $y$ is arbitrary.
\end{lemma}

\begin{proof}
The total array $R_s=P_s+Q_s$ of any vector satisfies
$\|(R_{si})_i\|_p\le R_s$, by Minkowski.  Consequently its $p$th
powers sum to at most $R_s^p$ on any antichain below $s$.
This follows first for finite antichains by iterating the child
inequality, and then for countable ones by increasing finite subsets.

There is a feasible pair for $y$ which uses its minimal arrays off
$A$ and whose two values at $s\in A$ both equal
\[
 \left(\sum_{i\in H:\,s\preceq i}
                  (P_i(y)+Q_i(y))^p\right)^{1/p}.
\]
Grouping boundary descendants by the first child proves the child
inequalities.  The antichain estimate just proved makes the sums
finite.  Minimality, \eqref{var:boundary-signs}, and Minkowski give
\[
 Y\le2\|(P_i(y)+Q_i(y))_{i\in H}\|_p
   \le2\bigl(\|(v_i^+)_{i\in H}\|_p+
                \|(v_i^-)_{i\in H}\|_p\bigr).
\]
Choose a sign $\sigma$ with $\|(v_i^\sigma)_{i\in H}\|_p\ge Y/4$.
For this sign, put
\[
 R_s=P_s(x+\sigma y)+Q_s(x+\sigma y),\qquad
 d_s=d_s^\sigma\ (s\in A),\qquad d_i=v_i^\sigma\ (i\in H).
\]
Let $L_s$ be the old value at $s\in A$ of its chosen tight array.
Subtract its old equality from the new child inequality.  For
children in $A$, use $(a+d)^p-a^p\ge d^p$, and for boundary children
use that the old value was zero.  It follows that
\[
 \sum_{i\text{ child of }s}d_i^p
 \le (L_s+d_s)^p-L_s^p
 \le pR_s^{p-1}d_s.
\]
Here and below a node in $A$ with no boundary descendants makes no
contribution.  Define $z_s=d_s/R_s$ when $R_s>0$, and $z_s=0$
otherwise.  All child values vanish when $R_s=0$; for positive $R_s$
we have
\begin{equation}\label{var:normalized-recursion}
 z_s\ge\frac1p\sum_{i\text{ child of }s}
             \left(\frac{R_i}{R_s}\right)^p z_i^p.
\end{equation}
The weights in this sum have total at most one.

To keep track of the boundary mass, define for $s\in A\cup H$
\[
 W_s=\frac{1}{R_s^p}
       \sum_{i\in H:\,s\preceq i}(v_i^\sigma)^p
 \quad(R_s>0),
 \qquad W_s=0\quad(R_s=0).
\]
The antichain bound gives $0\le W_s\le1$.  At $s\in H$ we have
$W_s=z_s^p$, and at $s\in A$ the quantity $W_s$ is the weighted sum
of the child quantities with the weights in
\eqref{var:normalized-recursion}.
For a node with remaining height $m=h-|s|$, we claim
\begin{equation}\label{var:normalized-iteration}
 z_s\ge C_m W_s^{p^m},\qquad
 C_0=1,\qquad C_m=p^{-1}C_{m-1}^p.
\end{equation}
It holds at boundary nodes because $z_s=W_s^{1/p}$,
$0\le W_s\le1$, and $C_m\le1$.  At nodes of $A$ of depth $h$,
$W_s=0$.
For the induction step, insert the child bounds into
\eqref{var:normalized-recursion} and apply Jensen to the power
$p^m$, placing any missing weight at zero.  This proves
\eqref{var:normalized-iteration}, including countably many children
by the same Jensen inequality for a probability measure on a
countable set.

At the root,
\[
 C_h=p^{-(1+p+\cdots+p^{h-1})}\ge\exp(-(\mathrm e-1)),
 \qquad p^{h+1}\le2\mathrm e<6.
\]
Indeed $p^h\le\mathrm e$ and $\log p\le p-1$ bound the first
geometric sum after taking logarithms.  If $Y=0$, the conclusion is
immediate from $M_\pm\ge B$.  Otherwise let $M=M_\sigma=R_o$.
The selected sign gives $W_o\ge(Y/(4M))^p$, so $Y/(4M)\le1$.
Equation \eqref{var:normalized-iteration} therefore yields
\[
 M-B=2d_o\ge
 2M\exp(-(\mathrm e-1))\left(\frac{Y}{4M}\right)^6.
\]
Since $M\le U$,
\[
 M^2-B^2\ge M(M-B)\ge c_6M^{-4}Y^6
                         \ge c_6U^{-4}Y^6.
\]
The other squared endpoint norm is at least $B^2$, proving
\eqref{var:sixth-component-eq}.  If $x=0$, then $U^2=2Y^2$,
and the same inequality holds directly because $c_6\le8$.
\end{proof}

We can now sum the component estimate without a constant depending
on the number of components or their heights.

\begin{theorem}\label{var:sixth-stability}
Let $x\in X_{\mathrm v}$ have finite total support.  In each component
where $x^{(h)}\ne0$, choose a finite ancestral set $A_h$ containing
that support.  Suppose $y\in X_{\mathrm v}$ vanishes on every such
$A_h$.  With
$S=\|x+y\|_{\mathrm v}^2+\|x-y\|_{\mathrm v}^2$, one has
\begin{equation}\label{var:sixth-sum}
 S-2\|x\|_{\mathrm v}^2
 \ge c_6 S^{-2}\|y\|_{\mathrm v}^6,
 \qquad c_6=2\exp(-(\mathrm e-1))4^{-6}.
\end{equation}
The right side is zero when $S=0$.
\end{theorem}

\begin{proof}
Set $Y_h=\nu_h(y^{(h)})$ and
$U_h^2=\nu_h(x^{(h)}+y^{(h)})^2+
        \nu_h(x^{(h)}-y^{(h)})^2$.
Lemma~\ref{var:sixth-component} applies also to components with
$x^{(h)}=0$.  Summing it reduces the assertion to
\[
 \sum_{U_h>0}\frac{Y_h^6}{U_h^4}
 \ge S^{-2}\left(\sum_hY_h^2\right)^3.
\]
For $S>0$, this is Jensen for the cube with weights $U_h^2/S$ and
values $Y_h^2/U_h^2$.  Components with $U_h=0$ have both vectors
zero and can be omitted.  The case $S=0$ is immediate.
\end{proof}

\subsection{Accumulating the exit energy}

The preceding proof chooses one sign and repeatedly raises a
normalized boundary quantity to a power.  We now keep both signs
and retain the mass at every boundary level after it has left the
head.  This yields a different recurrence and the cubic estimate.

\begin{lemma}\label{var:cubic-component}
Let $x,y\in E_h$, with $x$ supported on a finite ancestral set $A$
and $y$ zero on $A$.  Put
$M_\sigma=\nu_h(x+\sigma y)$, $B=\nu_h(x)$, and
$M=\max\{M_+,M_-\}$.  If $M>0$, then
\begin{equation}\label{var:cubic-component-eq}
 \frac{M_++M_-}{2}-B
 \ge c_3\frac{\nu_h(y)^3}{M^2},\qquad c_3=(6\mathrm e)^{-3}.
\end{equation}
This also holds for $x=0$ and arbitrary $y$.  When $M=0$, both
vectors are zero.
\end{lemma}

\begin{proof}
By homogeneity, assume $M=1$, so $\nu_h(y)\le1$.  Use $p=p_h$,
the common increments \eqref{var:common-increments}, and the same
choice of tight label at each node as above.  Denote that label by
$S$ in a calculation at its chosen node; the children in that
calculation use the parent's label.  For a nonterminal $s\in A$,
the subtraction that gave \eqref{var:normalized-recursion} gives
\begin{equation}\label{var:layer-increment}
 \sum_{i:\,si\in A}(d_{si}^\sigma)^p
 +\sum_{i:\,si\notin A} S_{si}(x+\sigma y)^p
 \le pS_s(x+\sigma y)^{p-1}d_s^\sigma.
\end{equation}
At an outside child, the two signs interchange $P,Q$.  Thus the sum
of the second term over signs is the sum of the $p$th powers of
both arrays for $y$, independently of the parent's label.

For $0\le j\le h$, and for $0\le j<h$ respectively, define
\[
 L_j^p=\sum_{\substack{s\in A\\|s|=j}}\sum_{\sigma=\pm}
                   (d_s^\sigma)^p,
 \qquad
 H_j^p=\sum_{\substack{s\in A\\|s|=j}}
             \sum_{i:\,si\notin A}
                  \bigl(P_{si}(y)^p+Q_{si}(y)^p\bigr).
\]
An empty sum is zero.  Summing \eqref{var:layer-increment} gives
\begin{equation}\label{var:layer-sum}
 L_{j+1}^p+H_j^p
 \le p\sum_{\substack{s\in A\\|s|=j}}\sum_{\sigma=\pm}
                    S_s(x+\sigma y)^{p-1}d_s^\sigma.
\end{equation}
The sum of $S_s(x+\sigma y)^p$ on the right is at most two:
for each sign, the sum of the $p$th powers of both arrays on a
level is at most $P_o(x+\sigma y)^p+Q_o(x+\sigma y)^p\le1$.
Moreover, the first-exit children form an antichain, so
\begin{equation}\label{var:exit-mass}
 \sum_{j<h}H_j^p
 \le P_o(y)^p+Q_o(y)^p\le\nu_h(y)^p\le1.
\end{equation}
The antichain argument includes infinitely many children by
increasing finite subsets, as in Lemma~\ref{var:sixth-component}.

We retain the exit mass from all preceding levels by setting
\[
 K_j=\left(L_j^p+\sum_{i<j}H_i^p\right)^{1/p},
 \qquad 0\le j\le h.
\]
Add $\sum_{i<j}H_i^p$ to \eqref{var:layer-sum}, bound it by
$p\sum_{i<j}H_i^{p-1}H_i$, and apply H\"older to this list of
terms together with the terms indexed by $(s,\sigma)$.
The $p/(p-1)$ powers of the first factors sum to at most
$2+1=3$, by the level bound and \eqref{var:exit-mass}; the
$p$th powers of the second factors sum to $K_j^p$.  Therefore
\begin{equation}\label{var:accumulated-recursion}
 K_{j+1}^p\le p\,3^{(p-1)/p}K_j.
\end{equation}
Iterating gives
\begin{equation}\label{var:accumulated-bound}
 K_h\le
 \bigl(p\,3^{(p-1)/p}\bigr)^{\sum_{i=1}^{h}p^{-i}}
 K_0^{p^{-h}}
 \le3\mathrm e\,K_0^{p^{-h}}.
\end{equation}
For the last inequality, use
$\sum_{i=1}^{h}p^{-i}\le1/(p-1)$ and
$p^{1/(p-1)}\le\mathrm e$.

It remains to relate the accumulated exit mass to the norm of $y$.
On $A$, the two arrays for $y$ are equal, since $y_s=0$.
Their common $p$th power at a nonterminal node $s\in A$ is
\[
 \sum_{i:\,si\in A}P_{si}(y)^p+
 \max\left\{\sum_{i:\,si\notin A}P_{si}(y)^p,
             \sum_{i:\,si\notin A}Q_{si}(y)^p\right\}.
\]
At nodes of $A$ of depth $h$ it is zero.  Bounding each maximum by the sum
of its two entries and iterating upward yields
$P_o(y)^p\le\sum_{j<h}H_j^p$.  Consequently
\[
 \nu_h(y)=2P_o(y)\le2K_h.
\]
Combining with \eqref{var:accumulated-bound}, and using
$p^h\le\mathrm e<3$ and $0\le\nu_h(y)/(6\mathrm e)\le1$, gives
\[
 K_0\ge\left(\frac{\nu_h(y)}{6\mathrm e}\right)^{p^h}
      \ge\left(\frac{\nu_h(y)}{6\mathrm e}\right)^3.
\]
Finally,
\[
 \frac{M_++M_-}{2}-B=d_o^++d_o^-\ge
 \bigl((d_o^+)^p+(d_o^-)^p\bigr)^{1/p}=K_0.
\]
This proves the normalized estimate and hence
\eqref{var:cubic-component-eq} by scaling.  For $x=0$ its left
side is $\nu_h(y)=M$, and the right side is $c_3M\le M$.
\end{proof}

\begin{theorem}\label{var:cubic-stability}
Let $x\in X_{\mathrm v}$ have finite total support.  In every
component with $x^{(h)}\ne0$, choose a finite ancestral set $A_h$
containing its support, and let $y\in X_{\mathrm v}$ vanish on these
sets.  If $R>0$ and $\|x+y\|_{\mathrm v},\|x-y\|_{\mathrm v}\le R$,
then
\begin{equation}\label{var:cubic-sum}
 R^2\ge\|x\|_{\mathrm v}^2+
       \frac{c_3}{2\sqrt2\,R}\|y\|_{\mathrm v}^3,
 \qquad c_3=(6\mathrm e)^{-3}.
\end{equation}
\end{theorem}

\begin{proof}
In component $h$ write
$B_h=\nu_h(x^{(h)})$, $Y_h=\nu_h(y^{(h)})$,
$M_{h,\sigma}=\nu_h(x^{(h)}+\sigma y^{(h)})$,
$M_h=\max_\sigma M_{h,\sigma}$, and
$m_h=(M_{h,+}+M_{h,-})/2$.
Lemma~\ref{var:cubic-component}, including its zero-center case,
and $m_h\ge M_h/2$ give
\[
 \frac{M_{h,+}^2+M_{h,-}^2}{2}-B_h^2
 \ge(m_h-B_h)(m_h+B_h)
 \ge\frac{c_3}{2}\frac{Y_h^3}{M_h}
\]
whenever $M_h>0$.  When $M_h=0$, both component vectors vanish.
H\"older gives
\[
 \sum_hY_h^2
 \le\left(\sum_{M_h>0}\frac{Y_h^3}{M_h}\right)^{2/3}
       \left(\sum_hM_h^2\right)^{1/3}.
\]
Since $\sum_hM_h^2\le2R^2$, summing the component inequalities
proves \eqref{var:cubic-sum}.  The same calculation for finite
partial sums and monotone convergence justifies all infinite sums.
\end{proof}

The finite-head hypotheses in these two theorems do not require
the perturbation to have finite support.  Their proofs already use
the completed component spaces and countable sums.  The following
consequence removes the finite-support restriction on the center
when the perturbations form a weakly null sequence. It bounds the norm
of the center from the size of weakly null midpoint perturbations.

\begin{proposition}\label{var:cubic-weak-null}
Let $x\in X_{\mathrm v}$ and $0<\varepsilon\le1$.  Suppose
$(y_j)$ is weakly null,
$\|y_j\|_{\mathrm v}\ge\varepsilon$, and
$\|x+y_j\|_{\mathrm v},\|x-y_j\|_{\mathrm v}\le1$ for every $j$.
Then
\begin{equation}\label{var:theta}
 \|x\|_{\mathrm v}\le\theta(\varepsilon)
 :=\left(1-\frac{c_3}{2\sqrt2}\varepsilon^3\right)^{1/2}<1.
\end{equation}
\end{proposition}

\begin{proof}
Fix $\eta>0$ and choose a vector $x'$ of finite total support with
$\|x-x'\|_{\mathrm v}<\eta$.  In its nonzero components, take
finite ancestral sets containing the support, and let $P$ be the
projection onto their union.  Proposition~\ref{var:space} makes
$P$ bounded and finite rank, so weak nullity gives
$\|Py_j\|_{\mathrm v}\to0$.  Set $y'_j=y_j-Py_j$.
The pair $x',y'_j$ satisfies the hypotheses of
Theorem~\ref{var:cubic-stability}, with
$R_j=1+\eta+\|Py_j\|_{\mathrm v}$.  Moreover,
$\liminf_j\|y'_j\|_{\mathrm v}\ge\varepsilon$.
First let $j\to\infty$ in \eqref{var:cubic-sum} to obtain
\[
 (1+\eta)^2\ge\|x'\|_{\mathrm v}^2+
          \frac{c_3\varepsilon^3}{2\sqrt2(1+\eta)}.
\]
Then let $\eta\downarrow0$.  This gives \eqref{var:theta}.
\end{proof}

\paragraph{The averaged midpoint modulus.}
The finite-head estimate also gives
\[
 \widehat\delta_{\|\cdot\|_{\mathrm v}}(t)
 \ge\frac{c_3}{24\sqrt2}\,t^3\qquad(0<t<1).
\]
To see this, fix a finite unit center $x$ and let $F$ be the kernel of
the finite coordinate projection onto ancestral sets containing its
support in its nonzero components. For $y\in F$ with
$\|y\|_{\mathrm v}=1$, put
$R=\max\{\|x+ty\|_{\mathrm v},\|x-ty\|_{\mathrm v}\}$.
The component inequalities $M_\sigma\ge B$ in
\eqref{var:common-increments}, together with the zero-center
components, give $\|x\pm ty\|_{\mathrm v}\ge1$. Also
$R\le1+t<2$. Theorem~\ref{var:cubic-stability} therefore yields
\[
 R(R+1)(R-1)\ge\frac{c_3}{2\sqrt2}\,t^3,
 \qquad R-1\ge\frac{c_3}{12\sqrt2}\,t^3.
\]
Since the other endpoint norm is at least one, their mean is at least
$1+(R-1)/2$. For an arbitrary unit center $x_0$, approximate it by a
finite unit center $x$. Replacing $x$ by $x_0$ loses at most
$\|x-x_0\|_{\mathrm v}$ in the mean, uniformly over unit $y\in F$.
The subspace $F$ is closed and finite codimensional by
Proposition~\ref{var:space}. Letting this approximation error tend to
zero in the definition of the modulus proves the claim; unit
directions suffice by the radius monotonicity established in the
introduction.
